\subsection{Sequences - Examples}
\tsEx{Explicit \& Recursive Sequences}{
    Explicit: $a_n=\frac{1}{n+1}$ gives $1,\frac12,\frac13,\frac14,\dots\to0$.
    \tsPoint Recursive: $a_0=1,\ a_{n+1}=\frac12 a_n \Rightarrow a_n=\left(\frac12\right)^n\to0$.
}
\tsEx{Oscillating vs.\ Divergent Sequences}{
    $a_n=3+\frac{(-1)^n}{n}$ oscillates around $3$ with shrinking amplitude, so $\lim_{n\to\infty}a_n=3$ (converges).
    \tsPoint $a_n=n$: for every $M>0$ choose $N>M$, then $n>N\Rightarrow a_n>M$, so $\lim_{n\to\infty}a_n=\infty$ (diverges).
}
\tsEx{Subsequence \& Accumulation Points}{
    $a_n=\frac1n$ with even indices $n_k=2k$ gives the subsequence $a_{n_k}=\frac1{2k}\to0$, same limit as $(a_n)$.
    \tsPoint $a_n=(-1)^n$ diverges but has two accumulation points $\pm1$, via $a_{2k}=1$ and $a_{2k+1}=-1$.
}
\tsEx{Limit Laws \& Sandwich Theorem}{
    $a_n=\frac1n\to0,\ b_n=2+\frac1n\to2 \Rightarrow a_n+b_n\to2,\ a_nb_n\to0$.
    \tsPoint $b_n=\frac{\sin n}{n}$: since $-\frac1n\le b_n\le\frac1n\to0$, the sandwich theorem gives $b_n\to0$.
}
\tsEx{Monotone Convergence}{
    $a_n=1-\frac1n$ satisfies $0\le a_n<1$ and $a_{n+1}>a_n$: bounded and monotone increasing, hence convergent with $\lim_{n\to\infty}a_n=1$.
}
